\section{Description du problème}
\label{sec:problem}

La section~\ref{sec:introduction} a établi que les modèles sémantiques de bâtiments libèrent une valeur significative
mais sont coûteux à créer.
Cette section examine le problème des coûts en termes concrets.
Nous décrivons ce que sont les points BACnet, quantifions le travail requis pour les mapper, et identifions le fossé de recherche qui a empêché une solution automatisée.

\subsection{Le problème du mappage des points BACnet}
\label{subsec:bacnet-problem}

Un Système d'automatisation du bâtiment (GTB/BAS) expose les données opérationnelles à travers une hiérarchie de \emph{points},
des capteurs individuels, des actionneurs, des consignes, des alarmes et des indicateurs d'état.
Dans BACnet (le protocole ouvert dominant pour les contrôles des bâtiments commerciaux), chaque point est un objet typé
(Entrées analogiques, Sorties binaires, Valeurs analogiques, etc.)\ avec une adresse unique sur le réseau.
Un immeuble de bureaux commercial typique possède entre 500 et 5{,}000 points BACnet ; un grand hôpital ou
une installation de campus peut en avoir des dizaines de milliers~\cite{Trenbath2022}.

Le \emph{problème de mappage} est le suivant : une liste brute de points BACnet ne contient aucune information standard sur
l'équipement physique auquel chaque point appartient, le rôle qu'il joue dans le système, ou comment cet équipement
se connecte à d'autres composants.
Un point nommé AHU-1.SA-T indique à un ingénieur humain \og probablement la température de l'air soufflé de
l'Unité de traitement d'air 1 \fg{}, mais seulement si l'ingénieur connaît déjà la convention de dénomination utilisée par ce
particular installateur GTB.
Il est courant de trouver des conventions qui diffèrent entre les entrepreneurs, entre les projets, et même entre
les systèmes au sein d'un même bâtiment.

Reconstituer le tableau complet du bâtiment --- quels points appartiennent à quel équipement, ce que signifie chaque
point, et comment l'équipement se connecte en un système CVC complet --- nécessite la lecture des plans CVC tels que
construits, l'interrogation de la programmation GTB et l'application d'un jugement d'ingénierie expert.
Ce processus est manuel, chronophage et dépendant de l'expertise du domaine.

Modbus, un protocole encore plus ancien et plus simple utilisé dans les contrôleurs CVC hérités, les compteurs et le matériel industriel,
expose les données sous forme de registres numérotés sans aucune métadonnée descriptive, rendant
le défi de mappage encore plus difficile que BACnet.

\subsection{Le coût du mappage manuel}
\label{subsec:cost}

Le mappage manuel des points BACnet et la mise en service du GTB sont largement documentés comme des facteurs de coût majeurs dans
les projets de construction et de rénovation de bâtiments.
Le Tableau~\ref{tab:commissioning-costs} résume les chiffres clés issus de la recherche industrielle et des
études de cas publiées.

\begin{table}[htbp]
\centering
\caption{Données clés de coûts et d'efforts pour la mise en service manuelle des systèmes GTB et le mappage des points.}
\label{tab:commissioning-costs}
\small
\begin{tabular}{p{5.5cm}p{3.5cm}p{4cm}}
\toprule
\textbf{Poste de coût} & \textbf{Chiffres} & \textbf{Source} \\
\midrule
Mise en service GTB en \% de la construction & 0.5\%--1.5\% du total ; 8\%--15\% de la valeur du projet GTB
  & \cite{BCA2024} \\
\addlinespace
Étude de cas : école de 150\,000\,pi\textsuperscript{2}
  & \$225{,}000 de mise en service (\$1.50/pi\textsuperscript{2})
  & \cite{PingCx2024} \\
\addlinespace
Installation GTB complète moyenne
  & \$7.76/pi\textsuperscript{2} (plage : \$5.20--\$14.18)
  & \cite{Trenbath2022} \\
\addlinespace
Surcoût d'intégration de GTB existant
  & +20\%--40\% par rapport au coût de base
  & \cite{Envigilance2023} \\
\addlinespace
Appels de service supplémentaires liés à un mauvais mappage
  & 12--15 appels supplémentaires/an ; +30\%--45\% de temps de dépannage
  & \cite{PingCx2024} \\
\addlinespace
Dérive de performance énergétique due à une mauvaise mise en service
  & 15\%--30\% de gaspillage énergétique dans les 3 premières années ; \$0.20--\$0.50/pi\textsuperscript{2}/an
  & \cite{PingCx2024} \\
\addlinespace
Dépenses imprévues la première année (étude de cas école)
  & \$175{,}000, presque égal au coût de mise en service initial
  & \cite{PingCx2024} \\
\addlinespace
Pénurie de techniciens CVC (2024--2025)
  & 110{,}000 postes non pourvus
  & \cite{AccessCoins2025} \\
\bottomrule
\end{tabular}
\end{table}

Pour illustrer la nature cumulée de ces coûts : l'école de 150{,}000\,pi\textsuperscript{2} dans l'étude de cas
a payé \$225{,}000 pour la mise en service au départ, puis a engagé environ \$175{,}000 en dépenses imprévues
au cours de la première année en raison d'une documentation insuffisante des points, doublant presque le coût effectif de la
phase de mise en service~\cite{PingCx2024}.
Selon le NREL~\cite{Trenbath2022}, l'installation GTB complète coûte en moyenne \$7.76 par pied carré
aux États-Unis, avec la prime d'intégration des systèmes hérités ajoutant encore 20 à 40
pour cent sur les projets de rénovation~\cite{Envigilance2023} (données de 2022 et 2023).

Les approches manuelles sont également incomplètes par conception.
La pratique industrielle teste généralement seulement 10 à 20 pour cent des unités terminales lors de la mise en service~\cite{PingCx2024} ;
chaque réduction de 10 pour cent de la couverture des tests est corrélée à une augmentation d'environ 6 à 8 pour cent
des problèmes post-occupation.
Le résultat est un problème omniprésent et chronique : les bâtiments commencent leur vie opérationnelle avec
des systèmes de contrôle sous-documentés et partiellement vérifiés, et le coût de correction augmente avec le temps.
Wang et al.~\cite{Wang2017} décrivent le mappage manuel des points comme \og laborieux et coûteux, présentant
un obstacle majeur aux innovations dans le contrôle des bâtiments \fg{}, validant l'ampleur du problème même
là où les chiffres précis dépendent de l'entrepreneur.

\subsection{Le fossé de recherche et les solutions émergentes}
\label{subsec:research-gap}

Malgré l'ampleur et l'impact économique du problème de mappage manuel, il n'existait pas de solution automatisée
de bout en bout largement déployée au moment de la revue systématique de 2017 de Wang et al.~\cite{Wang2017}, qui a constaté
que le domaine manquait :
\begin{itemize}[leftmargin=*, itemsep=2pt]
  \item d'une formulation complète du problème avec des approches de solution intégrées,
  \item de jeux de données de test standardisés et de métriques de performance, et
  \item de méthodes bien alignées avec les exigences de mappage du monde réel.
\end{itemize}

La vague d'IA à partir de 2023 a catalysé une nouvelle génération de tentatives pour combler ce fossé.
Le Tableau~\ref{tab:landscape} passe en revue les efforts commerciaux et académiques les plus pertinents et positionne
SI-Mapper dans le paysage.
Les orientations publiées par la Building Commissioning Association~\cite{BCA2024} confirment que
la vérification point à point reste une étape obligatoire sur chaque nouveau projet GTB, soulignant que le
problème persiste même lorsque des solutions partielles émergent.

\begin{table}[htbp]
\centering
\caption{Paysage concurrentiel pour le mappage automatisé des points GTB et la génération de modèles sémantiques (2024--2026).}
\label{tab:landscape}
\small
\begin{tabular}{p{3cm}p{3cm}p{3cm}p{4cm}}
\toprule
\textbf{Solution} & \textbf{Catégorie} & \textbf{Entrée} & \textbf{Écart par rapport à SI-Mapper} \\
\midrule

PingCx~\cite{PingCx2024}
  & Mise en service fonctionnelle automatisée et tests point à point
  & Connexion BACnet en direct
  & Teste si les points répondent correctement ; ne construit \emph{pas} de modèle sémantique ni n'identifie la topologie des équipements \\
\addlinespace

BrickLLM~\cite{Perini2025}
  & Génération d'ontologie Brick pilotée par l'IA
  & Description du bâtiment en texte libre
  & Nécessite une saisie en langage naturel précise ; ne traite pas les exports BACnet CSV, les plans CVC, ni ne produit du 223P \\
\addlinespace

BuildingMOTIF~\cite{BuildingMOTIF_NREL}
  & SDK open-source du DOE pour la création et validation de modèles 223P / Brick / Haystack
  & RDF rédigé par des experts ; listes de points BACnet (manuel)
  & Boîte à outils de développeur, pas un moteur d'inférence par IA ; nécessite un expert du domaine pour piloter la création du modèle \\
\addlinespace

BIM2RDF / Open223~\cite{Open223}
  & Conversion BIM vers 223P
  & Fichier BIM Revit / IFC
  & Nécessite un BIM complet et précis en entrée ; ne traite pas les bâtiments où seuls les plans tels que construits et les CSV BACnet existent \\
\addlinespace

Willow~\cite{Willow2025}
  & Plateforme de jumeau numérique IA avec classification ML des équipements
  & Flux GTB / IoT en direct
  & Utilise l'ontologie propriétaire DTDL, pas 223P ; nécessite une connexion de données en direct et une intégration guidée par un humain ; plateforme commerciale financée à \$110M \\
\addlinespace

SkySpark / Project Haystack~\cite{SkySpark}
  & Plateforme d'analytique avec étiquetage sémantique
  & Étiquettes Haystack appliquées par un ingénieur
  & L'étiquetage est manuel ; fournit un moteur d'analytique une fois étiqueté, pas un outil de mappage automatisé \\
\addlinespace

\textbf{SI-Mapper (ce travail)}
  & \textbf{Pipeline IA de bout en bout pour la génération de modèles 223P}
  & \textbf{Exports CSV de points BACnet + plans CVC tels que construits}
  & \textbf{Aucun système comparable} ne traite le pipeline complet, des exports BACnet bruts et des plans d'ingénierie jusqu'à un graphe ASHRAE 223P validé \\

\bottomrule
\end{tabular}
\end{table}

Le tableau révèle un fossé structurel clair.
Les solutions existantes soit s'arrêtent avant la modélisation sémantique (PingCx automatise les tests mais ne produit aucune
ontologie), exigent des entrées propres qui n'existent pas en pratique (BrickLLM a besoin d'un texte précis ;
BIM2RDF a besoin d'un fichier BIM complet), ou sont des boîtes à outils de développement destinées aux experts plutôt que de
l'automatisation clé en main (BuildingMOTIF).
Les plateformes commerciales comme Willow abordent le problème dans la direction opposée, de haut en bas,
en commençant par les données en direct dans de grands portefeuilles, en utilisant une ontologie propriétaire qui n'est pas alignée avec
ASHRAE 223P ou les normes d'interopérabilité émergentes requises pour l'intégration des centrales électriques virtuelles.

Les programmes des laboratoires nationaux (NREL, LBNL, PNNL, NIST) ont réalisé des progrès significatifs dans la définition
de la norme 223P et des outils associés à travers BuildingMOTIF~\cite{BuildingMOTIF_NREL,DOE_BuildingMOTIF},
mais le cadrage explicite du DOE reconnaît que l'objectif est toujours de \og réduire les coûts et le temps d'installation \fg{},
une admission que l'outillage actuel n'atteint pas encore cela sans une implication substantielle d'experts.

Le fossé que SI-Mapper comble est donc spécifique et étroit :
\emph{génération automatisée de bout en bout de modèles sémantiques ASHRAE 223P à partir des artefacts qui
existent réellement sur le terrain}, exports CSV BACnet et plans d'ingénierie tels que construits, sans
nécessiter un fichier BIM préexistant, des descriptions en langage naturel propres, ou un sémanticien expert.

\subsection{Pourquoi cela importe pour le réseau électrique}
\label{subsec:grid-relevance}

Le problème du mappage des points n'est pas seulement un défi opérationnel des bâtiments.
C'est une barrière directe à la transition énergétique.

\textbf{Les alternatives non filaires} et les \textbf{centrales électriques virtuelles} agrègent les ressources énergétiques
distribuées, notamment les charges CVC des bâtiments commerciaux, pour fournir des services de réseau électrique équivalents à
la construction de nouvelles centrales électriques ou de lignes de transport.
Les systèmes CVC consomment environ 40 pour cent de l'énergie totale des bâtiments commerciaux et peuvent être
modulés pour répondre aux signaux du réseau, déplaçant la charge loin des périodes de pointe de demande.
Globalement, cette flexibilité représente une capacité de réseau électrique potentielle énorme.

Le marché des centrales électriques virtuelles croît rapidement.
Les programmes nord-américains de centrales électriques virtuelles totalisaient 37.5\,GW de capacité flexible derrière le compteur
en 2025, soit environ 14 pour cent de plus qu'en 2024, avec des déploiements actifs en hausse de plus de 33
pour cent dans la même période~\cite{WoodMac2025}.
Le soutien politique s'accélère : dix législatures d'États ont introduit des projets de loi sur les centrales électriques virtuelles en
2024~\cite{RMI_VP3_2025,SEPA_NCCETC2025} ; les services publics du Massachusetts mettent en œuvre un
Fonds de compensation des services de réseau de \$50\,millions pour promouvoir les alternatives non filaires~\cite{RMI_VP3_2025} ;
et la Commission de l'énergie de Californie a accordé une subvention EPRI de \$2\,millions spécifiquement pour démontrer
les centrales électriques virtuelles pour les bâtiments commerciaux incluant les contrôles CVC~\cite{CEC_EPRI2024}.

\textbf{Comment les agrégateurs fonctionnent aujourd'hui.}
Deux types d'agrégateurs participent à ces programmes.
Les \emph{agrégateurs matériels} fabriquent des appareils contrôlables (thermostats intelligents, batteries, chargeurs VE)
et agrègent leur propre parc installé en un portefeuille de ressources distribuables.
Les \emph{agrégateurs logiciels} agissent comme des proxys de protocole, connectant des appareils hétérogènes de
plusieurs fournisseurs sous une couche de gestion unique.
Les deux types vendent des services de capacité et de réponse à la demande aux services publics et aux opérateurs de réseau, qui les envoient
en utilisant des signaux standardisés tels qu'OpenADR ou IEEE~2030.5.
Ce modèle fonctionne efficacement pour les actifs résidentiels et les flottes d'appareils identiques car
l'agrégateur contrôle le matériel d'extrémité et connaît déjà ses capacités.

\textbf{Les bâtiments commerciaux sont structurellement différents.}
Le système CVC d'un bâtiment commercial n'est pas un appareil que l'agrégateur a fabriqué.
Il a été conçu par un bureau d'études, installé par un entrepreneur en génie mécanique, et programmé par un
fournisseur GTB (Siemens, Johnson Controls, ABB, Niagara, ou équivalent).
Chacune de ces parties a un intérêt légitime à protéger la fiabilité du système : le fournisseur GTB
garantit le temps de fonctionnement ; l'ingénieur de mise en service a conçu les séquences de contrôle ; l'opérateur du bâtiment
est légalement responsable du confort des occupants.
Aucun d'entre eux n'autorisera volontiers une partie externe à écrire des consignes arbitraires dans leur système
sans un accord formel et une frontière technique claire.
Il en résulte que chaque inscription d'un bâtiment commercial aujourd'hui nécessite un projet de négociation et d'intégration
personnalisé, bâtiment par bâtiment \cite{DOE_VPP2025}.
Le rapport Liftoff des centrales électriques virtuelles du DOE identifie qu'il n'existe aucune standardisation industrielle.
Les fournisseurs de centrales électriques virtuelles ont construit leurs solutions de zéro, puis
ont travaillé au cas par cas avec les services publics pour développer des solutions personnalisées, résultant en des systèmes rigides
qui ne peuvent pas être facilement répliqués dans de nouveaux domaines~\cite{DOE_VPP2025}.

\textbf{Comment une ontologie ASHRAE 223P change la logistique.}
Un modèle sémantique ASHRAE 223P d'un bâtiment n'est pas un remplacement du système de contrôle, c'est un
contrat lisible par machine qui décrit ce qui est contrôlable, comment y accéder, et dans quelles
conditions.
Considérez la chaîne opérationnelle qu'il permet :

\begin{enumerate}[leftmargin=*, itemsep=2pt]
  \item \textbf{Découverte des actifs.} L'ontologie déclare chaque composant contrôlable (ventilateurs, bobines,
        boîtes VAV, registres) comme une entité typée avec sa position dans la topologie CVC.
        Un agrégateur peut l'interroger pour répondre à \og quelles charges ce bâtiment contient et quelles sont
        leurs capacités nominales ? \fg{} sans visite sur site ni révision manuelle de la documentation.
  \item \textbf{Accès à la télémétrie en direct.} Le mécanisme de \emph{référence externe} d'ASHRAE 223P relie
        chaque propriété sémantique directement à son adresse d'objet BACnet.
        Un capteur de température de l'air soufflé dans le modèle porte l'adresse du point BACnet physique
        qui signale sa valeur.
        L'agrégateur lit les données de capteurs en direct en suivant ces références, sans code d'intégration
        personnalisé par bâtiment.
  \item \textbf{Estimation de la flexibilité.} Avec la topologie et la télémétrie en direct disponibles, un
        algorithme peut calculer combien de charge le bâtiment peut délester, pendant combien de temps, et avec quel
        décalage thermique, avant que les contraintes de confort ne soient violées.
        Cette estimation alimente directement l'optimisation de la répartition de l'agrégateur.
  \item \textbf{Répartition contrôlée.} L'ontologie identifie également les adresses BACnet des
        consignes contrôlables (consigne de température de l'air soufflé, commande de vitesse du ventilateur, consigne de refroidissement).
        Lorsque l'opérateur de réseau envoie un signal de délestage OpenADR, l'agrégateur le traduit
        en écritures BACnet ciblées précisément aux adresses que spécifie l'ontologie, dans
        les limites convenues contractuellement avec le propriétaire du bâtiment.
        Le GTB continue d'exécuter sa logique de fiabilité ; l'agrégateur ajuste uniquement les consignes
        qui ont été explicitement autorisées dans le cadre de l'accord.
\end{enumerate}

